\documentclass[11pt,a4paper]{article}

\usepackage[T1]{fontenc}
\usepackage[utf8]{inputenc}
\usepackage{lmodern}
\usepackage[a4paper,margin=1in]{geometry}
\usepackage{amsmath,amssymb,mathtools}
\usepackage{booktabs}
\usepackage{enumitem}
\usepackage{hyperref}
\usepackage{setspace}

\hypersetup{
  colorlinks=true,
  linkcolor=blue,
  urlcolor=blue,
  citecolor=blue
}

\setstretch{1.08}
\setlength{\parskip}{0.45em}
\setlength{\parindent}{0pt}

\title{Bragg Interferometer Simulation Formula Reference}
\author{}
\date{}

\begin{document}

\maketitle
\tableofcontents
\newpage

\section{Scope of the Model}

This document describes the effective model implemented by the GUI-based Bragg atom-interferometer simulator used in this project. The model is designed for:

\begin{itemize}[leftmargin=1.5em]
  \item fringe prediction,
  \item noise budgeting,
  \item sensitivity studies,
  \item parameter scans over pulse quality and technical noise.
\end{itemize}

It is not a full momentum-lattice propagation model. Instead, it keeps only an effective two-port interferometer subspace and treats imperfect Bragg diffraction through effective transfer probabilities, loss terms, and diffraction phases.
The general Bragg-interferometer setting, as well as the appearance of parasitic ports, losses, and diffraction phases in nonideal Bragg diffraction, are well documented in the literature on Bragg atom interferometry and diffraction phases \cite{Giltner1995,Muller2008,Estey2015,Parker2016,Jenewein2022}.

\section{Simulation Workflow}

The calculation proceeds in the following order:

\begin{enumerate}[leftmargin=1.8em]
  \item Define the interferometer geometry and the pulse sequence parameters.
  \item Build the laser-frequency-noise and mirror-acceleration-noise spectra.
  \item Convert these technical noise spectra into interferometer phase-noise spectra.
  \item Integrate the phase-noise spectra through the interferometer transfer function to obtain RMS phase noise.
  \item Build one effective complex matrix for each Bragg pulse.
  \item Propagate the initial state through the $\pi/2-\pi-\pi/2$ sequence.
  \item Compute upper-port and lower-port populations for each Monte Carlo shot.
  \item Repeat the calculation over a scanned optical phase and average the shots.
  \item Fit the resulting fringe with a cosine model.
\end{enumerate}

Each step is written explicitly below.

\section{Interferometer Parameters}

The interferometer-level parameters are:

\begin{itemize}[leftmargin=1.5em]
  \item $n$: Bragg order.
  \item $\lambda$: laser wavelength.
  \item $T$: pulse separation between the first and second pulse, and between the second and third pulse.
  \item $a_{\mathrm{eff}}$: effective acceleration projected along the interferometer axis.
  \item $\phi_{\mathrm{offset}}$: additional static phase offset.
  \item $\phi_0$: scanned optical phase.
\end{itemize}

These are entered in the GUI as:

\begin{itemize}[leftmargin=1.5em]
  \item Bragg order,
  \item Wavelength,
  \item Pulse separation $T$,
  \item Effective acceleration,
  \item Additional phase offset.
\end{itemize}

\section{Effective Wavevector}

The single-photon wavevector magnitude is
\begin{equation}
  k = \frac{2\pi}{\lambda}.
\end{equation}

The effective Bragg interferometer wavevector is
\begin{equation}
  k_{\mathrm{eff}} = 2nk = 2n\frac{2\pi}{\lambda}.
\end{equation}
This is the standard large-momentum-transfer scaling for Bragg diffraction, where the interferometric phase is enhanced by the diffraction order $n$ \cite{Giltner1995,Muller2008}.

\textbf{Parameter meanings}
\begin{itemize}[leftmargin=1.5em]
  \item $k$: single-photon wavevector magnitude.
  \item $k_{\mathrm{eff}}$: effective two-photon Bragg wavevector used in the interferometer phase.
  \item $n$: Bragg order.
  \item $\lambda$: wavelength.
\end{itemize}

\section{Deterministic Interferometer Phase}

The effective deterministic phase used in the model is
\begin{equation}
  \phi_{\mathrm{base}} = k_{\mathrm{eff}} a_{\mathrm{eff}} T^2 + \phi_{\mathrm{offset}}.
\end{equation}

The scanned optical phase enters as
\begin{equation}
  \phi_{\mathrm{scan}} = n\phi_0.
\end{equation}

Therefore the deterministic phase before shot-to-shot technical noise is
\begin{equation}
  \phi_{\mathrm{det}} = \phi_{\mathrm{base}} + \phi_{\mathrm{scan}}.
\end{equation}
The three-pulse Mach-Zehnder phase structure and the laser-phase imprinting formalism follow the standard light-pulse interferometer treatment \cite{Cheinet2008}. The explicit $n$ scaling is the Bragg large-momentum-transfer generalization discussed in Bragg-interferometer references such as Refs.~\cite{Giltner1995,Muller2008}.

\textbf{Parameter meanings}
\begin{itemize}[leftmargin=1.5em]
  \item $\phi_{\mathrm{base}}$: inertial plus static phase.
  \item $a_{\mathrm{eff}}$: effective acceleration.
  \item $T$: pulse separation.
  \item $\phi_{\mathrm{offset}}$: additional user-defined phase offset.
  \item $\phi_0$: scanned optical phase.
  \item $n$: Bragg order.
\end{itemize}

\section{Noise Spectra Input Convention}

The simulator accepts noise spectra in power spectral density form.

For laser-frequency noise, the PSD unit is
\begin{equation}
  \mathrm{Hz}^2/\mathrm{Hz}.
\end{equation}

For mirror acceleration noise, the PSD unit is
\begin{equation}
  (\mathrm{m/s}^2)^2/\mathrm{Hz}.
\end{equation}

All dB inputs use the convention
\begin{equation}
  \mathrm{PSD}_{\mathrm{dB}} = 10\log_{10}\left(\frac{S}{S_{\mathrm{ref}}}\right),
\end{equation}
with:
\begin{itemize}[leftmargin=1.5em]
  \item $S_{\mathrm{ref}} = 1~\mathrm{Hz}^2/\mathrm{Hz}$ for laser-frequency noise,
  \item $S_{\mathrm{ref}} = 1~(\mathrm{m/s}^2)^2/\mathrm{Hz}$ for mirror acceleration noise.
\end{itemize}

The inverse conversion is
\begin{equation}
  S = 10^{\mathrm{PSD}_{\mathrm{dB}}/10}.
\end{equation}

\textbf{Parameter meanings}
\begin{itemize}[leftmargin=1.5em]
  \item $S$: linear PSD.
  \item $\mathrm{PSD}_{\mathrm{dB}}$: PSD written in dB.
  \item $S_{\mathrm{ref}}$: reference PSD unit used in the dB convention.
\end{itemize}

\section{Background Noise Model}

The simulator decomposes the technical noise background into:
\begin{itemize}[leftmargin=1.5em]
  \item a white term,
  \item a flicker term,
  \item a random-walk term,
  \item a sum of Lorentzian peaks.
\end{itemize}

The full PSD is modeled as
\begin{equation}
  S(f) = S_0 + \frac{C_1}{f} + \frac{C_2}{f^2} + \sum_i S_{\mathrm{peak},i}(f).
\end{equation}

\textbf{Parameter meanings}
\begin{itemize}[leftmargin=1.5em]
  \item $f$: Fourier frequency.
  \item $S_0$: white PSD level.
  \item $C_1$: flicker PSD coefficient referenced at $1~\mathrm{Hz}$.
  \item $C_2$: random-walk PSD coefficient referenced at $1~\mathrm{Hz}$.
  \item $S_{\mathrm{peak},i}(f)$: PSD contribution of peak $i$.
\end{itemize}

\subsection{White Term}

The white-noise contribution is
\begin{equation}
  S_{\mathrm{white}}(f) = S_0.
\end{equation}

\textbf{Parameter meanings}
\begin{itemize}[leftmargin=1.5em]
  \item $S_0$: constant PSD floor.
\end{itemize}

\subsection{Flicker Term}

The flicker-noise contribution is
\begin{equation}
  S_{\mathrm{flicker}}(f) = \frac{C_1}{f}.
\end{equation}

\textbf{Parameter meanings}
\begin{itemize}[leftmargin=1.5em]
  \item $C_1$: PSD coefficient equal to the flicker PSD at $f=1~\mathrm{Hz}$.
  \item $f$: Fourier frequency.
\end{itemize}

\subsection{Random-Walk Term}

The random-walk contribution is
\begin{equation}
  S_{\mathrm{rw}}(f) = \frac{C_2}{f^2}.
\end{equation}

\textbf{Parameter meanings}
\begin{itemize}[leftmargin=1.5em]
  \item $C_2$: PSD coefficient equal to the random-walk PSD at $f=1~\mathrm{Hz}$.
  \item $f$: Fourier frequency.
\end{itemize}

\section{Lorentzian Peak Model}

Each narrow-band peak is modeled by a Lorentzian PSD:
\begin{equation}
  S_{\mathrm{peak}}(f) =
  \frac{S_{\mathrm{peak,center}}}
  {1 + \left(\frac{f-f_0}{\Gamma}\right)^2}.
\end{equation}

\textbf{Parameter meanings}
\begin{itemize}[leftmargin=1.5em]
  \item $S_{\mathrm{peak,center}}$: peak-center PSD in linear units.
  \item $f_0$: center frequency of the peak.
  \item $\Gamma$: width parameter of the Lorentzian.
  \item $f$: Fourier frequency.
\end{itemize}

In the GUI, $S_{\mathrm{peak,center}}$ is entered as $\mathrm{psd\_db}$, and the simulator converts it back to linear PSD through
\begin{equation}
  S_{\mathrm{peak,center}} = 10^{\mathrm{psd\_db}/10}.
\end{equation}

\section{Total ASD Used Internally}

The simulator computes the total linear PSD first:
\begin{equation}
  S_{\mathrm{total}}(f) =
  S_0 + \frac{C_1}{f} + \frac{C_2}{f^2} + \sum_i S_{\mathrm{peak},i}(f).
\end{equation}

The amplitude spectral density used internally is then
\begin{equation}
  \mathrm{ASD}_{\mathrm{total}}(f) = \sqrt{S_{\mathrm{total}}(f)}.
\end{equation}

\textbf{Parameter meanings}
\begin{itemize}[leftmargin=1.5em]
  \item $S_{\mathrm{total}}(f)$: total linear PSD.
  \item $\mathrm{ASD}_{\mathrm{total}}(f)$: total amplitude spectral density.
\end{itemize}

\section{Laser-Frequency Noise to Phase Noise}

Let $S_{\nu}(f)$ be the laser-frequency-noise PSD. The simulator constructs the interferometer phase-noise PSD as
\begin{equation}
  S_{\phi,\mathrm{laser}}(f) = \frac{n^2 S_{\nu}(f)}{f^2}.
\end{equation}
This phase-noise treatment is an effective frequency-domain reduction consistent with the three-pulse sensitivity-function formalism used for light-pulse atom interferometers \cite{Cheinet2008}.

\textbf{Parameter meanings}
\begin{itemize}[leftmargin=1.5em]
  \item $S_{\phi,\mathrm{laser}}(f)$: interferometer phase-noise PSD arising from laser frequency noise.
  \item $S_{\nu}(f)$: laser-frequency-noise PSD.
  \item $n$: Bragg order.
  \item $f$: Fourier frequency.
\end{itemize}

This is the effective model used in the code. It compresses the laser-noise transfer into a single phase-noise PSD before the final transfer-function weighting.

\section{Mirror Acceleration Noise to Phase Noise}

Let $S_a(f)$ be the mirror acceleration PSD. The mirror displacement PSD is implicitly modeled via
\begin{equation}
  S_z(f) = \frac{S_a(f)}{(2\pi f)^4}.
\end{equation}

The corresponding interferometer phase-noise PSD is
\begin{equation}
  S_{\phi,\mathrm{vib}}(f) = k_{\mathrm{eff}}^2 S_z(f)
  = \frac{k_{\mathrm{eff}}^2 S_a(f)}{(2\pi f)^4}.
\end{equation}
The conversion from acceleration noise to displacement noise and then to phase noise is again based on the standard sensitivity-function analysis of three-pulse atom interferometers \cite{Cheinet2008}.

\textbf{Parameter meanings}
\begin{itemize}[leftmargin=1.5em]
  \item $S_a(f)$: mirror acceleration PSD.
  \item $S_z(f)$: mirror displacement PSD.
  \item $S_{\phi,\mathrm{vib}}(f)$: interferometer phase-noise PSD from mirror vibration.
  \item $k_{\mathrm{eff}}$: effective wavevector.
  \item $f$: Fourier frequency.
\end{itemize}

\section{Interferometer Transfer Function}

For the effective three-pulse sequence, the simulator uses the frequency-domain weighting
\begin{equation}
  H(f) = \left|1 - 2e^{-i2\pi fT} + e^{-i4\pi fT}\right|.
\end{equation}

The weighted phase-noise integrand is therefore
\begin{equation}
  I(f) = |H(f)|^2 S_{\phi}(f).
\end{equation}

This is evaluated separately for laser noise and vibration noise.
This transfer-function description is directly motivated by the standard sensitivity-function formalism for three-pulse atom interferometers \cite{Cheinet2008}.

\textbf{Parameter meanings}
\begin{itemize}[leftmargin=1.5em]
  \item $H(f)$: magnitude of the interferometer transfer function.
  \item $I(f)$: weighted phase-noise integrand.
  \item $T$: pulse separation.
  \item $S_{\phi}(f)$: phase-noise PSD of one technical noise source.
\end{itemize}

\section{Integrated Phase Noise}

The simulator integrates the weighted phase-noise PSD over a logarithmic frequency grid:
\begin{equation}
  \sigma_{\phi}^2 = \int_{f_{\min}}^{f_{\max}} |H(f)|^2 S_{\phi}(f)\,df.
\end{equation}

The RMS phase noise is
\begin{equation}
  \sigma_{\phi} = \sqrt{\sigma_{\phi}^2}.
\end{equation}
The code evaluates this integral numerically on a logarithmic grid; the underlying formal relation follows Ref.~\cite{Cheinet2008}.

This is done independently for:
\begin{itemize}[leftmargin=1.5em]
  \item laser-frequency noise,
  \item mirror-vibration noise.
\end{itemize}

\textbf{Parameter meanings}
\begin{itemize}[leftmargin=1.5em]
  \item $\sigma_{\phi}$: RMS phase noise contributed by one noise source.
  \item $f_{\min}, f_{\max}$: integration bounds.
  \item $|H(f)|^2 S_{\phi}(f)$: transfer-function-weighted phase-noise PSD.
\end{itemize}

\section{Pulse-Level Effective Parameters}

Each Bragg pulse is parameterized by:
\begin{itemize}[leftmargin=1.5em]
  \item $P$: transfer probability,
  \item $L$: loss probability,
  \item $\phi_d$: diffraction phase,
  \item $\sigma_P$: shot-to-shot standard deviation of the transfer probability,
  \item $\sigma_{\phi_d}$: shot-to-shot standard deviation of the diffraction phase.
\end{itemize}

These are defined separately for:
\begin{itemize}[leftmargin=1.5em]
  \item the first beam splitter,
  \item the mirror pulse,
  \item the second beam splitter.
\end{itemize}
The quantities $P$, $L$, and $\phi_d$ should be interpreted as effective parameters. They are not copied from a single exact multistate Bragg Hamiltonian, but rather serve as a reduced description of the transfer efficiency, parasitic-port loss, and diffraction phase discussed in the Bragg-diffraction literature \cite{Muller2008,Estey2015,Parker2016,Jenewein2022}.

\section{Shot-to-Shot Pulse Fluctuations}

For each Monte Carlo shot, the transfer probability is sampled as
\begin{equation}
  P' \sim \mathcal{N}(P,\sigma_P),
\end{equation}
and clipped to the interval $[0,1]$.

The diffraction phase is sampled as
\begin{equation}
  \phi_d' \sim \mathcal{N}(\phi_d,\sigma_{\phi_d}).
\end{equation}

\textbf{Parameter meanings}
\begin{itemize}[leftmargin=1.5em]
  \item $P'$: shot-specific transfer probability.
  \item $\phi_d'$: shot-specific diffraction phase.
  \item $P$: mean transfer probability.
  \item $\phi_d$: mean diffraction phase.
  \item $\sigma_P$: standard deviation of the transfer probability.
  \item $\sigma_{\phi_d}$: standard deviation of the diffraction phase.
\end{itemize}

\section{Effective Pulse Matrix}

Each pulse is represented by the effective complex matrix
\begin{equation}
  U(P',\phi_d',L) =
  \sqrt{1-L}
  \begin{bmatrix}
    \sqrt{1-P'} & -i e^{-i\phi_d'} \sqrt{P'} \\
    -i e^{i\phi_d'} \sqrt{P'} & \sqrt{1-P'}
  \end{bmatrix}.
\end{equation}
This matrix should be read as a phenomenological two-port model inspired by the population loss, multiport output, and diffraction-phase effects analyzed in Refs.~\cite{Muller2008,Estey2015,Parker2016,Jenewein2022}. It is not presented as an exact equation reproduced from any one of those papers.

\textbf{Parameter meanings}
\begin{itemize}[leftmargin=1.5em]
  \item $U$: effective pulse matrix.
  \item $P'$: shot-specific transfer probability.
  \item $L$: loss probability.
  \item $\phi_d'$: shot-specific diffraction phase.
\end{itemize}

Interpretation:
\begin{itemize}[leftmargin=1.5em]
  \item diagonal elements describe the amplitude that remains in the same effective port,
  \item off-diagonal elements describe transfer between the two effective ports,
  \item the prefactor $\sqrt{1-L}$ removes population that leaks into untracked parasitic momentum states.
\end{itemize}

\section{Initial State}

The interferometer starts in the lower basis state of the effective two-port model:
\begin{equation}
  \psi_0 =
  \begin{bmatrix}
    1 \\
    0
  \end{bmatrix}.
\end{equation}

\textbf{Parameter meanings}
\begin{itemize}[leftmargin=1.5em]
  \item $\psi_0$: initial state vector before the first pulse.
\end{itemize}

\section{Total Shot Phase}

For each Monte Carlo shot, the total interferometer phase applied to the final pulse is
\begin{equation}
  \phi_{\mathrm{shot}} =
  \phi_{\mathrm{det}}
  + \delta\phi_{\mathrm{laser}}
  + \delta\phi_{\mathrm{vib}}.
\end{equation}

with
\begin{equation}
  \delta\phi_{\mathrm{laser}} \sim \mathcal{N}(0,\sigma_{\phi,\mathrm{laser}})
\end{equation}
and
\begin{equation}
  \delta\phi_{\mathrm{vib}} \sim \mathcal{N}(0,\sigma_{\phi,\mathrm{vib}}).
\end{equation}

\textbf{Parameter meanings}
\begin{itemize}[leftmargin=1.5em]
  \item $\phi_{\mathrm{shot}}$: total shot-specific interferometer phase.
  \item $\phi_{\mathrm{det}}$: deterministic phase from inertial and user-defined terms.
  \item $\delta\phi_{\mathrm{laser}}$: laser-noise phase draw.
  \item $\delta\phi_{\mathrm{vib}}$: vibration-noise phase draw.
  \item $\sigma_{\phi,\mathrm{laser}}$: RMS laser-noise phase.
  \item $\sigma_{\phi,\mathrm{vib}}$: RMS vibration-noise phase.
\end{itemize}

\section{Three-Pulse State Propagation}

The three effective pulse matrices are:
\begin{equation}
  U_1 = U(P_1',\phi_{d,1}',L_1),
\end{equation}
\begin{equation}
  U_2 = U(P_2',\phi_{d,2}',L_2),
\end{equation}
\begin{equation}
  U_3 = U(P_3',\phi_{d,3}' + \phi_{\mathrm{shot}},L_3).
\end{equation}

The final state is
\begin{equation}
  \psi_f = U_3 U_2 U_1 \psi_0.
\end{equation}
The sequence structure is the standard three-pulse Mach-Zehnder order used in light-pulse atom interferometry \cite{Cheinet2008}, while the specific reduced matrices are the effective Bragg model introduced in this document.

\textbf{Parameter meanings}
\begin{itemize}[leftmargin=1.5em]
  \item $U_1, U_2, U_3$: first beam splitter, mirror, and second beam splitter matrices.
  \item $\phi_{\mathrm{shot}}$: shot-specific interferometer phase added to the last pulse.
  \item $\psi_f$: final state vector after the full sequence.
\end{itemize}

\section{Output-Port Populations}

If
\begin{equation}
  \psi_f =
  \begin{bmatrix}
    \psi_A \\
    \psi_B
  \end{bmatrix},
\end{equation}
then the absolute tracked-port populations are
\begin{equation}
  P_A = |\psi_A|^2,
\end{equation}
\begin{equation}
  P_B = |\psi_B|^2.
\end{equation}

The effective loss channel is
\begin{equation}
  P_{\mathrm{loss}} = 1 - (P_A + P_B).
\end{equation}

\textbf{Parameter meanings}
\begin{itemize}[leftmargin=1.5em]
  \item $P_A$: lower effective output-port population.
  \item $P_B$: upper effective output-port population.
  \item $P_{\mathrm{loss}}$: population leaked out of the tracked two-port model.
\end{itemize}

\section{Normalised Port Populations}

The plotted upper-port normalised population is
\begin{equation}
  P_{B,\mathrm{norm}} = \frac{P_B}{P_A + P_B}.
\end{equation}

The normalised lower-port population is
\begin{equation}
  P_{A,\mathrm{norm}} = \frac{P_A}{P_A + P_B}.
\end{equation}

Since the denominator is the same, the two satisfy
\begin{equation}
  P_{A,\mathrm{norm}} = 1 - P_{B,\mathrm{norm}}.
\end{equation}

This identity is exact inside the model, even when $P_{\mathrm{loss}}$ is nonzero.

\textbf{Parameter meanings}
\begin{itemize}[leftmargin=1.5em]
  \item $P_{B,\mathrm{norm}}$: normalised upper-port population.
  \item $P_{A,\mathrm{norm}}$: normalised lower-port population.
\end{itemize}

\section{Monte Carlo Averaging Over Shots}

At each scanned phase point, the simulator runs $N_{\mathrm{shot}}$ independent shots and computes:
\begin{equation}
  \overline{P}_{B,\mathrm{norm}} =
  \frac{1}{N_{\mathrm{shot}}}
  \sum_{j=1}^{N_{\mathrm{shot}}} P_{B,\mathrm{norm}}^{(j)},
\end{equation}
\begin{equation}
  \overline{P}_{A,\mathrm{norm}} =
  \frac{1}{N_{\mathrm{shot}}}
  \sum_{j=1}^{N_{\mathrm{shot}}} P_{A,\mathrm{norm}}^{(j)}.
\end{equation}

The absolute means are similarly defined:
\begin{equation}
  \overline{P}_A =
  \frac{1}{N_{\mathrm{shot}}}
  \sum_{j=1}^{N_{\mathrm{shot}}} P_A^{(j)},
\end{equation}
\begin{equation}
  \overline{P}_B =
  \frac{1}{N_{\mathrm{shot}}}
  \sum_{j=1}^{N_{\mathrm{shot}}} P_B^{(j)}.
\end{equation}

\textbf{Parameter meanings}
\begin{itemize}[leftmargin=1.5em]
  \item $N_{\mathrm{shot}}$: number of Monte Carlo shots per phase point.
  \item the overbar denotes a Monte Carlo mean over shots.
\end{itemize}

\section{Shot Scatter}

The plotted shot scatter is the sample standard deviation of the normalised port population at a fixed scanned phase:
\begin{equation}
  \sigma_{\mathrm{scatter}} =
  \sqrt{
    \frac{1}{N_{\mathrm{shot}}-1}
    \sum_{j=1}^{N_{\mathrm{shot}}}
    \left(P_{\mathrm{norm}}^{(j)} - \overline{P}_{\mathrm{norm}}\right)^2
  }.
\end{equation}

Because
\begin{equation}
  P_{A,\mathrm{norm}} = 1 - P_{B,\mathrm{norm}},
\end{equation}
the normalised upper-port and lower-port scatters are equal in this model.

\textbf{Parameter meanings}
\begin{itemize}[leftmargin=1.5em]
  \item $\sigma_{\mathrm{scatter}}$: sample standard deviation at one scanned phase point.
  \item $P_{\mathrm{norm}}^{(j)}$: normalised port population in shot $j$.
  \item $\overline{P}_{\mathrm{norm}}$: mean normalised port population.
\end{itemize}

\section{Cosine Fringe Fit}

After averaging over shots, the simulator fits the displayed normalised fringe with
\begin{equation}
  P_{\mathrm{fit}}(\phi_0) = C_0 + C\cos(n\phi_0 + \phi_{\mathrm{fit}}).
\end{equation}
The cosine fringe structure is the standard output form of a two-path light-pulse atom interferometer. In a Bragg interferometer the scanned optical phase is multiplied by the diffraction order $n$ through the large-momentum-transfer phase scaling \cite{Giltner1995,Muller2008}.

\textbf{Parameter meanings}
\begin{itemize}[leftmargin=1.5em]
  \item $C_0$: fitted offset.
  \item $C$: fitted fringe amplitude.
  \item $n$: Bragg order.
  \item $\phi_0$: scanned optical phase.
  \item $\phi_{\mathrm{fit}}$: fitted phase offset.
\end{itemize}

The reported contrast is defined as the peak-to-peak excursion of the displayed fringe over the scanned phase range:
\begin{equation}
  \mathcal{C}_{\mathrm{pp}} = P_{\max} - P_{\min}.
\end{equation}

\textbf{Parameter meanings}
\begin{itemize}[leftmargin=1.5em]
  \item $\mathcal{C}_{\mathrm{pp}}$: reported peak-to-peak contrast.
  \item $P_{\max}$: maximum value of the displayed normalised fringe over the scan.
  \item $P_{\min}$: minimum value of the displayed normalised fringe over the scan.
\end{itemize}

For an ideal cosine sampled densely over a full period, this satisfies
\begin{equation}
  \mathcal{C}_{\mathrm{pp}} = 2C,
\end{equation}
where $C$ is the cosine-fit amplitude defined above.

\section{Lower-Port Versus Upper-Port Display}

The upper-port display uses the curve
\begin{equation}
  P_{\mathrm{display}} = P_{B,\mathrm{norm}}.
\end{equation}

The lower-port display uses the curve
\begin{equation}
  P_{\mathrm{display}} = P_{A,\mathrm{norm}} = 1 - P_{B,\mathrm{norm}}.
\end{equation}

Therefore:
\begin{itemize}[leftmargin=1.5em]
  \item the two normalised port curves are complementary,
  \item they have the same scatter magnitude,
  \item their fitted cosine amplitudes have equal magnitude,
  \item their fitted phase offsets differ by $\pi$ modulo $2\pi$.
\end{itemize}

\section{Effective Rabi-Model Mapping}

When pulse parameters are constructed from an effective two-state Rabi model, the transfer probability is estimated as
\begin{equation}
  P =
  \frac{\Omega_{\mathrm{eff}}^2}
  {\Omega_{\mathrm{eff}}^2 + \Delta^2}
  \sin^2\left(
    \frac{\tau}{2}
    \sqrt{\Omega_{\mathrm{eff}}^2 + \Delta^2}
  \right).
\end{equation}
This is the standard off-resonant two-level Rabi-transition result. In the Bragg context, effective Rabi frequencies, losses, and associated phase shifts are discussed in detail in Ref.~\cite{Muller2008}.

\textbf{Parameter meanings}
\begin{itemize}[leftmargin=1.5em]
  \item $\Omega_{\mathrm{eff}}$: effective Rabi frequency.
  \item $\Delta$: detuning.
  \item $\tau$: pulse duration.
  \item $P$: effective transfer probability used in the pulse matrix.
\end{itemize}

This mapping is provided as a convenience function. The main GUI workflow accepts $P$ directly.

\section{Numerical Integration Grid}

The noise integration uses a logarithmic frequency grid:
\begin{equation}
  f_k \in [f_{\min}, f_{\max}],
\end{equation}
with $N_f$ points spaced logarithmically between the bounds.

\textbf{Parameter meanings}
\begin{itemize}[leftmargin=1.5em]
  \item $f_{\min}$: lower frequency cutoff.
  \item $f_{\max}$: upper frequency cutoff.
  \item $N_f$: number of frequency-grid points.
\end{itemize}

\section{Approximations and Limits}

The formulas above correspond to an effective model with the following approximations:
\begin{enumerate}[leftmargin=1.8em]
  \item Only two effective interferometer ports are propagated explicitly.
  \item Additional Bragg momentum ports are absorbed into effective loss.
  \item Phase noise is reduced to Gaussian shot-to-shot phase draws after transfer-function integration.
  \item Pulse imperfections are represented by effective transfer probability and effective diffraction phase.
  \item The reported shot scatter is Monte Carlo scatter of the model outputs, not a full atom-detection noise model.
\end{enumerate}

These approximations make the simulator fast and interpretable, but they should not be confused with a full multistate Bragg propagation treatment.

\begin{thebibliography}{99}

\bibitem{Cheinet2008}
P.~Cheinet, B.~Canuel, F.~Pereira Dos Santos, A.~Gauguet, F.~Leduc, and A.~Landragin,
``Measurement of the sensitivity function in a time-domain atomic interferometer,''
\textit{IEEE Transactions on Instrumentation and Measurement} \textbf{57}(6), 1141--1148 (2008).
DOI: \href{https://doi.org/10.1109/TIM.2007.915148}{10.1109/TIM.2007.915148}.

\bibitem{Giltner1995}
D.~M.~Giltner, R.~W.~McGowan, and S.~A.~Lee,
``Atom Interferometer Based on Bragg Scattering from Standing Light Waves,''
\textit{Physical Review Letters} \textbf{75}, 2638--2641 (1995).
DOI: \href{https://doi.org/10.1103/PhysRevLett.75.2638}{10.1103/PhysRevLett.75.2638}.

\bibitem{Muller2008}
H.~M\"uller, S.-w.~Chiow, and S.~Chu,
``Atom-wave diffraction between the Raman-Nath and the Bragg regime: Effective Rabi frequency, losses, and phase shifts,''
\textit{Physical Review A} \textbf{77}, 023609 (2008).
DOI: \href{https://doi.org/10.1103/PhysRevA.77.023609}{10.1103/PhysRevA.77.023609}.

\bibitem{Estey2015}
B.~Estey, C.~Yu, P.-C.~Kuan, and H.~M\"uller,
``High-Resolution Atom Interferometers with Suppressed Diffraction Phases,''
\textit{Physical Review Letters} \textbf{115}, 083002 (2015).
DOI: \href{https://doi.org/10.1103/PhysRevLett.115.083002}{10.1103/PhysRevLett.115.083002}.

\bibitem{Parker2016}
R.~H.~Parker, C.~Yu, B.~Estey, W.~Zhong, E.~Huang, and H.~M\"uller,
``Controlling the multiport nature of Bragg diffraction in atom interferometry,''
\textit{Physical Review A} \textbf{94}, 053618 (2016).
DOI: \href{https://doi.org/10.1103/PhysRevA.94.053618}{10.1103/PhysRevA.94.053618}.

\bibitem{Jenewein2022}
J.~Jenewein, S.~Hartmann, A.~Roura, and E.~Giese,
``Bragg-diffraction-induced imperfections of the signal in retroreflective atom interferometers,''
\textit{Physical Review A} \textbf{105}, 063316 (2022).
DOI: \href{https://doi.org/10.1103/PhysRevA.105.063316}{10.1103/PhysRevA.105.063316}.

\end{thebibliography}

\end{document}
